% Автор разбора: Павел Кротков

\begin{frame}
  \frametitle{Задача <<Различные числа>>}
  \begin{center}
    \includegraphics{differ.eps}
  \end{center}
\end{frame}

\begin{frame}
  \frametitle{Над задачей работали}
  \begin{itemize}
    \item Идея задачи: Виталий Аксенов
    \item Текст условия: Антон Банных
    \item Тесты, проверяющая программа и др: Сергей Мельников
    \item Решения: Сергей Мельников, Сергей Поромов, Андрей Комаров
    \item Текст разбора: Павел Кротков
  \end{itemize}
\end{frame}

\begin{frame}
  \frametitle{Формулировка задачи}
  \begin{itemize}
    \item Даны два целых числа: $a$ и $b$
    \item Требуется определить количество различных чисел среди следующих:
      \begin{itemize}
        \item $a + b$, $b + a$
        \item $a - b$, $b - a$
        \item $a * b$, $b * a$
        \item $\frac{a}{b}$, $\frac{b}{a}$
        \item $a ^ b$, $b ^ a$
      \end{itemize}
  \end{itemize}
\end{frame}

\begin{frame}
  \frametitle{Идея решения}
  \begin{itemize}
    \item Проверим все эти числа и посчитаем различные
    \item Если какое-то число не считается --- не будем его вычислять, мы сможем сравнить это число с другими и не зная его точного значения
  \end{itemize}
\end{frame}

\begin{frame}
  \frametitle{Что значит не вычисляется?}
  \begin{itemize}
    \item Модуль числа слишком большой
    \item Модуль числа слишком маленький
  \end{itemize}
\end{frame}

\begin{frame}
  \frametitle{Что делать с очень маленьким числом?}
  \begin{itemize}
    \item Очень маленьким(меньше $10^{-18}$) может получиться только число $a ^ b$ или $b ^ a$
    \item Если $a ^ b$ или $b ^ a$ очень мало и одно из чисел $a$ и $b$ меньше нуля, а другое больше, то $a ^ b$ и $b ^ a$ различны
    \item Если $a ^ b$ или $b ^ a$ очень мало и числа $a$ и $b$ меньше нуля, то $a ^ b = b ^ a$ только в случае равенства $a$ и $b$
  \end{itemize}
\end{frame}

\begin{frame}
  \frametitle{Что делать с очень большим числом?}
  \begin{itemize}
    \item Очень большим(больше $10^{18}$) может получиться только число $a ^ b$ или $b ^ a$
    \item Если $a ^ b$ или $b ^ a$ очень велико, то $a ^ b = b ^ a$ только в случае равенства $a$ и $b$
  \end{itemize}
\end{frame}

\begin{frame}
  \frametitle{Итого}
  \begin{itemize}
    \item Считаем значения, которые мы можем посчитать, и ищем количество различных.
    \item Если мы получили очень большие или маленькие числа, то скажем, сколько их
    \item Складываем ответы
  \end{itemize}
\end{frame}

\begin{frame}
  \frametitle{Спасибо за внимание}
  \begin{itemize}
    \item Вопросы?
  \end{itemize}
\end{frame}

